\section{平凡化映射}

出人意料的是，在纯规范理论中，
平凡化映射在一定程度上可以被显式地构造出来。
本节解释这一构造，假定规范作用量 $S(U)$
是 Wilson 圈（方块、矩形等）之和。


\subsection{平凡化流}

若流 (3.2) 的生成元 $Z_t(U)$ 满足
\begin{equation}
  \int_0^t\rmd s\,\sum_{x,\mu}
  \bigl\{\du^a_{x,\mu}[Z_s(U)]^a(x,\mu)\bigr\}_{U=U_s}
  =tS(U_t)+C_t,
\end{equation}
其中 $C_t$ 可以依赖 $t$ 但不依赖于场，
则相应的积分变换满足
\begin{equation}
  S(\transt(V))-\ln\det\transts(V)=(1-t)S(\transt(V))-C_t.
\end{equation}
特别地，此时 $t=1$ 处的变换就是一个平凡化映射。

方程 (4.1) 是对流的生成元的一个相当隐式的条件。
然而，对 $t$ 求导后，它化为一个更易处理的形式
\begin{equation}
  \sum_{x,\mu}\bigl\{
  \du^a_{x,\mu}[Z_t(U)]^a(x,\mu)
  -t\du^a_{x,\mu}S(U)[Z_t(U)]^a(x,\mu)\bigr\}=
  S(U)+\dot{C}_t,
\end{equation}
其中只涉及 $t$ 时刻的生成元。
注意，微分条件 (4.3)
与流方程 (3.2) 蕴含式~(4.1)，
即只需找到满足式~(4.3) 的生成元 $Z_t(U)$ 即可。


\subsection{平凡化流的存在性}

方程 (4.3) 是生成元 $Z_t(U)$ 的一个非齐次线性偏微分方程。
由于它是标量方程，
可以预期它有许多解。
下文中，解将取如下形式：
\begin{equation}
  [Z_t(U)]^a(x,\mu)=-\du^a_{x,\mu}\Sflow_t(U),
\end{equation}
其中作用量 $\Sflow_t(U)$ 待定。

把拟设 (4.4) 代入式~(4.3)，便得到 Laplace 方程
\begin{align}
  &\Lop\Sflow_t=S+\dot{C}_t,
  \\[1.0ex]
  &\Lop=
  \sum_{x,\mu}\bigl\{
  -\du^a_{x,\mu}\du^a_{x,\mu}
  +t\left(\du^a_{x,\mu}S\right)\du^a_{x,\mu}\bigr\}
\end{align}
（为简单起见，现在常省略宗量 $U$）。
算符 $\Lop$ 是椭圆型的，并且关于标量积
\begin{equation}
   (\phi,\psi)=\int\rmD[U]\,\rme^{-tS(U)}\phi(U)^*\psi(U)
\end{equation}
是对称的。
因此 $\Lop$ 具有完备的可微本征函数组，
以及没有聚点的纯离散谱
（例如参见文献~\cite{Gilkey}，第 1.6 节）。
此外，由于
\begin{equation}
   (\phi,\Lop\phi)=\sum_{x,\mu}\bigl(\du^a_{x,\mu}\phi,\du^a_{x,\mu}\phi
   \bigr)\geq0,
\end{equation}
函数 $\phi(U)=1$ 是 $\Lop$ 唯一的零模，
而所有其他本征函数的本征值与原点之间被严格正的谱隙隔开。

现在若这样选取 $C_t$ 使得
\begin{equation}
   \dot{C}_t=-(1,S)/(1,1),
\end{equation}
则零模分量便从式~(4.5) 的右端移除，
于是
\begin{equation}
   \Sflow_t=\Lop^{-1}(S+\dot{C}_t)
\end{equation}
是一个良好定义的表达式，且是该方程的解。
$\Sflow_t$ 关于 $t$ 与 $U$ 的可微性
本质上来自 $\Lop$ 的椭圆性（附录~E）。
至此我们给出了平凡化流存在性的一个构造性证明。


\subsection{$t$ 的幂级数展开}

解 (4.10) 是良好定义的，但仍相当隐式，
因为它涉及作用在规范场函数上的算符的逆。
在本小节及下一小节中，
将把解按 $t$ 的幂解析地展开。

当把级数
\begin{equation}
  \Sflow_t=\sum_{k=0}^{\infty}t^k\Sflow^{(k)}
\end{equation}
代入式~(4.5) 时，
匹配给定 $t$ 阶的项便得到递推关系
\begin{align}
  \Lap\Sflow^{(0)}&=S+\dot{C}^{(0)},
  \\[1.0ex]
  \Lap\Sflow^{(k)}&
  =-\sum_{x,\mu}\du^a_{x,\mu}S\,\du^a_{x,\mu}\Sflow^{(k-1)}
  +\dot{C}^{(k)},\qquad
  k=1,2,\ldots,
\end{align}
其中 $\Sflow^{(k)}$ 为作用量。
Laplacian $\Lap$ 与 $4+1$ 维格点规范理论中
哈密顿算符的色-电部分相重合。
特别地，其本征函数是链接变量的
$\Group$ 表示函数的乘积。
例如，Wilson 圈之和与 Wilson 圈之积
是 $\Lap$ 的本征函数，
或者可以容易地（代数地）分解为本征函数。

于是，该递推关系的解
\begin{align}
  \Sflow^{(0)}&=\Lap^{-1}S,
  \\[1.0ex]
  \Sflow^{(k)}&=-\Lap^{-1}
  \sum_{x,\mu}\du^a_{x,\mu}S\,\du^a_{x,\mu}
  \Sflow^{(k-1)},\qquad k=1,2,\ldots,
\end{align}
便以 Wilson 圈之和与 Wilson 圈之积的形式得到。
注意，如 4.2 小节所述，
常数函数是 $\Lap$ 唯一的零模。
由于 $\Lap$ 的最小非零本征值为 $\frac{4}{3}$，
故式~(4.14),(4.15) 的右端
除一个无关紧要的可加常数外被唯一确定。


\subsection{Wilson 理论中\/ $\Sflow^{(0)}$ 与\/ $\Sflow^{(1)}$ 的计算}

作为说明，现在对方块作用量~\cite{Wilson}
\begin{equation}
  \Sw(U)=-\frac{1}{6}\beta\,\Wl_0(U)
\end{equation}
（其中 $\beta=6/g_0^2$ 为逆规范耦合）
显式算出级数 (4.11) 的前两项。
利用完备性关系 (A.5) 作一简短计算，
可知领头项为
\begin{equation}
  \Sflow^{(0)}=-\frac{1}{32}\beta\Wl_0=\frac{3}{16}\Sw.
\end{equation}
精确到这一阶，并且在时间参数 $t$ 的重新标度之内，
Wilson 理论中的平凡化流因而与威尔逊流重合。

\begin{figure}[htbp]
  \centering
  \includegraphics[width=8.5cm]{images/loops.pdf}
  \caption{%
Wilson 理论中对 $\Sflow^{(1)}$ 有贡献的圈与圈对的类别。
圈 $5-7$
位于格点的单个方格上。
其余所有的圈占据两个方格，
这两个方格可以共面，也可以在三维空间中相互垂直。
}
  \label{fig:loops}
\end{figure}

下一阶要计算的表达式是
\begin{equation}
  \Sflow^{(1)}=
  -\frac{1}{192}\beta^2\Lap^{-1}\sum_{x,\mu}\du^a_{x,\mu}\Wl_0\,
  \du^a_{x,\mu}\Wl_0.
\end{equation}
此处可能出现的 Wilson 圈及 Wilson 圈之积，
来源于两个共享一条链接的方格圈的收缩。
于是总共有七类圈与圈对 $\mathcal{C}_i$，$i=1,\ldots,7$
需要考虑（见图~2）。
把相应 Wilson 圈的迹求和后，
每一类 $\mathcal{C}_i$ 定义了一个作用量
\begin{align}
  \Wl_i&=\sum_{C\in\mathcal{C}_i}\tr\bigl\{U(C)\bigr\}
  \quad\text{若}\quad i=1,2,5,
  \\[1.0ex]
  \Wl_i&=\sum_{\{C,C'\}\in\mathcal{C}_i}\tr\bigl\{U(C)\bigr\}\tr\bigl\{U(C')\bigr\}
  \quad\text{若}\quad i=3,4,6,7,
\end{align}
其中 $U(C)$ 表示沿圈 $C$ 的链接变量的有序乘积。
这些方程中的求和遍及圈在格点上所有可能的位置。
取向相反的圈被视为不同，
并且都被包含在求和中。

再次利用恒等式 (A.5)，经过一些代数运算得到
\begin{multline}
  \sum_{x,\mu}\du^a_{x,\mu}\Wl_0\,\du^a_{x,\mu}\Wl_0
  =\Wl_1-\Wl_2-\frac{1}{3}\Wl_3+\frac{1}{3}\Wl_4\\
  -2\Wl_5+\frac{2}{3}\Wl_6-\frac{4}{3}\Wl_7
\end{multline}
（相差一个可加常数）。
此外，
\begin{align}
  \Lap\Wl_1&=8\Wl_1,
  \\[0.6ex]
  \Lap\Wl_2&=\frac{31}{3}\Wl_2+\Wl_4,
  \\[0.6ex]
  \Lap\Wl_3&=11\Wl_3-\Wl_1,
  \\[0.6ex]
  \Lap\Wl_4&=\frac{31}{3}\Wl_4+\Wl_2,
  \\[0.6ex]
  \Lap\Wl_5&=\frac{28}{3}\Wl_5+4\Wl_6,
  \\[0.6ex]
  \Lap\Wl_6&=\frac{28}{3}\Wl_6+4\Wl_5,
  \\[0.6ex]
  \Lap\Wl_7&=12\Wl_7+\text{常数}.
\end{align}
在这些函数张成的子空间中，
算符 $\Lap$ 很容易求逆，
于是得到结果
\begin{multline}
  \Sflow^{(1)}=\frac{1}{192}\beta^2\Bigl\{
  -\frac{4}{33}\Wl_1
  +\frac{12}{119}\Wl_2
  +\frac{1}{33}\Wl_3
  -\frac{5}{119}\Wl_4
  +\frac{3}{10}\Wl_5\\
  -\frac{1}{5}\Wl_6
  +\frac{1}{9}\Wl_7\Bigr\}.
\end{multline}
注意，此式中数值系数的小，
在某种程度上被类 $\mathcal{C}_i$
中每个格点上的圈数所补偿
（对于 $i=1,\ldots,4$、$i=5,6$、$i=7$，
圈数分别等于 $120$、$12$ 和 $6$）。


\subsection{若干补充评论}

\noindent
(a)~{\itshape 更高阶}。
作用量 $\Sflow^{(k)}$（$k\geq2$）可以按照上一小节的步骤
代数地计算。
但由于必须考虑把方格圈与 $k-1$ 阶的圈作收缩所产生的所有圈，
所需计算量随 $k$ 迅速增长。

\smallskip
\noindent
(b)~{\itshape 定域性与收敛性}。级数 (4.11) 是局域项的展开，
这些项在格点上的足迹随阶数 $k$ 成比例增大。
于是在展开收敛的 $t$ 值处，
作用量 $\Sflow_t$ 也保证是定域的。

附录 E 中的范数估计给出级数收敛半径的一个下界，
但这个界相当差，并且在无穷体积极限下消失。
尽管如此，下列猜想似乎是合理的：
在此极限中级数具有非零收敛半径，
因为式~(4.10) 中算符 $\Lop$ 的逆
在 $t=0$ 的复邻域内很可能保持有界。
然而要证明这一点，
还需要计入 $\Lop$ 的定域性质的分析。

\smallskip
\noindent
(c)~{\itshape 展开的截断}。
如果把级数 (4.11) 中所有 $k\geq n$ 阶的项都略去，
便得到一个近似平凡化的流，
它满足式~(4.2) 直至一个 $t^{n+1}$ 量级的可加修正。
只要 $t$ 足够小使得该修正被强烈压低，
这种情况下的作用量至少会被部分抵消。

\smallskip
\noindent
(d)~{\itshape 光滑化性质}。
威尔逊流满足
\begin{equation}
  {\rmd\over\rmd t}\Sw(U_t)=-\frac{3}{16}\sum_{x,\mu}
  \bigl\{\du^a_{x,\mu}\Sw(U)
  \du^a_{x,\mu}\Sw(U)\bigr\}_{U=U_t}\leq 0
\end{equation}
因而随着 $t$ 增大单调降低 Wilson 作用量。
于是，本节为 Wilson 理论构造的平凡化流
在 $t$ 的领头阶上对规范场具有光滑化效应。
另一方面，若沿反向追踪该流，
规范场则会趋于变得更粗糙。

\smallskip
\noindent
(e)~{\itshape 拓扑荷扇区}。
在格点 QCD 中，拓扑（瞬子）扇区并非场流形单独的性质，
而是预期在接近连续极限时动力学地涌现。
因此，平凡化映射把这些扇区完全``拉直''
并不与场空间的拓扑性质相矛盾。

\smallskip
\noindent
(f)~{\itshape 重整化群}。
把 Wilson 理论中的平凡化映射 $U=\trans_1(V)$
与它在另一规范耦合值处的逆复合，
便得到一组变换，其对作用量的唯一效应只是耦合常数的平移。
然而，这些变换的定域性质并不清楚，
而且可能与 Wilson 式``块自旋''变换的定域性大不相同。
